% ══════════════════════════════════════════════════════════════════════════════
%  EXPERIMENT 1 — SOLUBILITY                                     (Lesson 3, Term 2 Week 1)
% ══════════════════════════════════════════════════════════════════════════════
\experimentdiv{Solubility}{Testing which substances dissolve in water \textbullet{} Keeping a test fair \textbullet{} Recording observations in a table}{Lesson 3 \textbullet\ Term 2, Week 1 \textbullet\ Lab}\label{sec:1}
\lessoncodes{VC2S8U06 \textbullet{} VC2S8I01, I02, I04, I08}

\vocabhead
\begin{ctab}
\begin{tabular}{|SW{4.2cm}|W{11.6cm}|}
  \hline\hd & \thead{Word} & \thead{Definition (copy from the slide)} \\ \hline
  & \scaf{\textbf{Dissolve}} & \scaf{When a solid mixes into a liquid so completely that it can no longer be seen.} \\ \hline
  & \textbf{Soluble} & \\ \hline
  & \textbf{Insoluble} & \\ \hline
  & \textbf{Solution} & \\ \hline
  & \textbf{Suspension} & \\ \hline
\end{tabular}
\end{ctab}

\begin{lipanel}
\inquiryli{Plan a fair test by keeping the amount of solid and water the same; record observations in a table.}
\understandli{Classify substances as soluble or insoluble in water.}
\sctitle
\begin{itemize}[leftmargin=1.4em, itemsep=0pt, topsep=1pt]
  \item I can predict which substances will dissolve and give a reason.
  \item I can name the variables I kept the same.
  \item I can classify each substance as soluble or insoluble from my observations.
\end{itemize}
\end{lipanel}

\begin{skillsfocus}
\focus{QP1}{Predict which substances will dissolve.}
\focus{PC2}{Keep the amounts of solid and water the same.}
\focus{PC6}{Observe each mixture closely.}
\focus{PA1}{Record observations in a table.}
\focus{CO1}{Write the conclusion in the third person.}
\end{skillsfocus}

% ─── TECHNIQUE ────────────────────────────────────────────────────────────────
\exppart{Technique spotlight}{Blue}

\begin{technique}{Telling soluble from insoluble}{PC6}
After stirring, hold the test tube up against a white card and look through it.
\begin{enumerate}[leftmargin=1.6em, itemsep=1pt, topsep=3pt]
  \item \textbf{Clear} (you can see through it, even if it is coloured) and no solid left: the substance is \textbf{soluble}. You have made a \textbf{solution}.
  \item \textbf{Cloudy}, and the solid settles to the bottom if you wait: the substance is \textbf{insoluble}. You have made a \textbf{suspension}.
  \item Clear, but some solid stays on the bottom: some of it dissolved. Record it as \textbf{partly soluble}.
\end{enumerate}
Clear does not mean colourless. Blue copper sulfate solution is clear; milky water is not.
\end{technique}

\begin{workedex}{Classifying from an observation}
\textbf{Observation:} \emph{After stirring, the liquid was cloudy brown. After one minute a layer of brown solid had settled on the bottom.}

\smallskip
\textbf{Classification:} insoluble. The liquid was cloudy and the solid settled out, so it is a suspension, not a solution.
\end{workedex}

\begin{wedo}{Classifying from an observation}
\emph{After stirring, the liquid was clear and pale yellow. No solid could be seen.} Classify the substance and give your reason.
\answerbox{3cm}
\end{wedo}

\begin{technique}{Keeping a test fair}{PC2}
In a fair test you change \textbf{one} thing (the independent variable), measure or observe \textbf{one} thing (the dependent variable), and keep \textbf{everything else} the same (the controlled variables). If you used more water for salt than for sand, you could not tell whether the salt dissolved because of what it is or because it had more water.
\end{technique}

\begin{youdobox}[Naming the variables]
\task Complete the table for today's experiment. \skr{PC2}
\begin{ctab}
\begin{tabular}{|SW{4.4cm}|W{10.2cm}|}
  \hline\hd & \thead{Variable} & \thead{In this experiment} \\ \hline
  & \scaf{Independent (what I change)} & \scaf{The substance being tested} \\ \hline
  & Dependent (what I observe) & \\ \hline
  & Controlled (what I keep the same) & \\ \hline
  & Controlled & \\ \hline
  & Controlled & \\ \hline
\end{tabular}
\end{ctab}
\end{youdobox}

% ─── PLAN AND SAFETY / METHOD ────────────────────────────────────────────────
\exppart{Plan, safety and method}{Teal}

\textbf{Aim:} To find out which common substances are soluble in water.

\task Predict whether each substance in the results table will dissolve in water. Write your prediction in the \emph{Prediction} column of Table 1.1, then explain one of your predictions here. \skr{QP1}
\answerbox{4cm}

\subsubsection*{Equipment}
8 test tubes in a rack \textbullet{} 10\,mL measuring cylinder \textbullet{} spatula \textbullet{} stirring rod \textbullet{} white card \textbullet{} small samples of salt, sugar, sand, chalk powder, flour, instant coffee, bicarbonate soda and cooking oil \textbullet{} wash bottle of water

\begin{center}
\begin{tikzpicture}[scale=0.8]
  \begin{scope}[shift={(0,0)}]\eqTestTube\fill[black!8] (0.02,0.3) arc[start angle=180,end angle=360,radius=0.28] -- (0.58,1.6) -- (0.02,1.6) -- cycle;\end{scope}
  \begin{scope}[shift={(2.2,0)}]\eqTestTube\fill[black!30] (0.02,0.3) arc[start angle=180,end angle=360,radius=0.28] -- (0.58,1.6) -- (0.02,1.6) -- cycle;\end{scope}
  \begin{scope}[shift={(4.4,0)}]\eqTestTube\fill[black!70] (0.02,0.3) arc[start angle=180,end angle=360,radius=0.28] -- (0.58,0.5) -- (0.02,0.5) -- cycle;\end{scope}
  \foreach \x in {0,2.2,4.4}{\draw[Ink, line width=0.8pt] (\x,1.6) -- (\x+0.6,1.6);}
  \node[eqtext, align=center] at (0.3,-0.45) {clear:\\soluble};
  \node[eqtext, align=center] at (2.5,-0.45) {cloudy:\\suspension};
  \node[eqtext, align=center] at (4.7,-0.45) {solid settled:\\insoluble};
\end{tikzpicture}
\expfigcap{What to look for after stirring.}
\end{center}

\begin{infobox}{Safety}
Wear safety glasses. Never taste anything in the lab, even substances from the kitchen. Wipe up spills of powder or oil straight away so that nobody slips.
\end{infobox}

\subsubsection*{Method}
\begin{enumerate}[itemsep=2pt]
  \item Label eight test tubes with the names of the substances.
  \item Measure 10\,mL of water into each test tube.
  \item Add one level spatula of the first substance to its test tube.
  \item Stir for 30 seconds, then hold the test tube against the white card and observe it.
  \item Wait one minute and observe again. Record both observations in Table 1.1.
  \item Repeat steps 3 to 5 for each substance, using a clean stirring rod each time.
\end{enumerate}

% ─── RESULTS ─────────────────────────────────────────────────────────────────
\exppart{Results}{Amber}

\subsubsection*{Table 1.1: Do these substances dissolve in water? \skr{PA1}}
\begin{ctab}
\begin{tabular}{|SW{2.6cm}|W{2.2cm}|W{3.8cm}|W{3.2cm}|W{2.6cm}|}
  \hline\hd & \thead{Substance} & \thead{Prediction} & \thead{After stirring} & \thead{After 1 minute} & \thead{Soluble or insoluble?} \\ \hline
  & \scaf{Salt} & \scaf{Will dissolve} & \scaf{Clear, colourless; no solid seen} & \scaf{No change} & \scaf{Soluble} \\ \hline
  & Sugar & & & & \\ \hline
  & Sand & & & & \\ \hline
  & Chalk powder & & & & \\ \hline
  & Flour & & & & \\ \hline
  & Instant coffee & & & & \\ \hline
  & Bicarbonate soda & & & & \\ \hline
  & Cooking oil & & & & \\ \hline
\end{tabular}
\end{ctab}

% ─── ANALYSIS ────────────────────────────────────────────────────────────────
\exppart{Analysis and evaluation}{Coral}

\task Sort the substances you tested into two lists. \skr{PA1}
\begin{ctab}
\begin{tabular}{|W{7.4cm}|W{7.4cm}|}
  \hline\thead{Soluble in water} & \thead{Insoluble in water} \\ \hline
  \rule{0pt}{3.2cm} & \\ \hline
\end{tabular}
\end{ctab}

\needspace{10\baselineskip}
\task How many of your predictions were correct? Choose one that was wrong and suggest why you predicted it wrongly.
\answerbox{4cm}

\needspace{10\baselineskip}
\task Instant coffee made a dark liquid. A student said, ``It's dark, so it can't have dissolved.'' Explain why the student is wrong.
\answerbox{4cm}

\needspace{10\baselineskip}
\task Explain why it was important to use the same volume of water and the same amount of each solid. \skr{PC2}
\answerbox{4cm}

\needspace{12\baselineskip}
\task Write a conclusion for this experiment. Answer the aim, give two examples, and write in the third person (no \emph{I} or \emph{we}). \skr{CO1}
\answerbox{5cm}

\begin{yr8box}
\begin{keepext}
\exttask Oil did not dissolve in water, but it does dissolve in other liquids, such as eucalyptus oil. Plan a fair test to find out which of three liquids dissolves the most substances. Name your independent, dependent and controlled variables.
\answerbox{6cm}
\end{keepext}
\end{yr8box}

% ─── SKILL CHECK ─────────────────────────────────────────────────────────────
\exppart{Skill check}{NavyBlue}

\begin{skillcheck}
\scrow{QP1}{I can predict which substances will dissolve and give a reason.}
\scrow{PC2}{I can name the variables I kept the same.}
\scrow{PC6}{I can tell a solution from a suspension by looking.}
\scrow{PA1}{I recorded every observation in my table.}
\scrow{CO1}{I can write a conclusion in the third person.}
\end{skillcheck}

\task Which skill will you focus on in Experiment 2? Explain what you will do differently.
\answerbox{5cm}
